%%
%% Test: Box Styles — v1.1
%% Stage: boxes refactor (stages 1-4)
%% 
%% Purpose: Verify box styles:
%%   - Three colors: blue (theorems), orange (examples), green (solutions)
%%   - No conflicting definitions with mb-theorem
%%   - matincream background
%%   - fonttitle=\normalsize
%%   - tcolorbox works with xepersian
%% 

\documentclass{matinbook}

\title{تست جعبه‌های رنگی — نسخه ۱.۱}
\author{تیم MatinBook}
\date{۱۴۰۵}

\booktitle{تست جعبه‌ها}

\begin{document}

\maketitle

\chapter{بررسی جعبه‌های رنگی}
\label{chap:test-boxes}

این سند، جعبه‌های رنگی \texttt{mb-boxes} را بررسی می‌کند.

%----------------------------------------
\section{خانواده قضیه (آبی)}
\label{sec:theorem-family}

\begin{theorem}[قضیه نمونه]
    برای هر عدد حقیقی $x$، داریم $x^2 \geq 0$.
    \label{thm:boxes-test}
\end{theorem}

\begin{lemma}[لم نمونه]
    اگر $a > 0$ و $b > 0$، آنگاه $a + b > 0$.
\end{lemma}

\begin{corollary}[نتیجه نمونه]
    اگر $a > 0$، آنگاه $a^2 > 0$.
\end{corollary}

\begin{proposition}[گزاره نمونه]
    اگر $n$ زوج باشد، $n^2$ نیز زوج است.
\end{proposition}

%----------------------------------------
\section{تعریف (آبی روشن)}
\label{sec:definition}

\begin{definition}[عدد اول]
    عدد طبیعی بزرگ‌تر از ۱ که تنها دو مقسوم‌علیه مثبت دارد،
    \emph{عدد اول} نامیده می‌شود.
\end{definition}

%----------------------------------------
\section{مثال و نکته (نارنجی)}
\label{sec:example-remark}

\begin{example}[اعداد اول کوچک]
    اعداد $2, 3, 5, 7, 11, 13$ نمونه‌هایی از اعداد اول هستند.
\end{example}

\begin{remark}[نکته]
    عدد ۱ نه اول است و نه مرکب.
\end{remark}

%----------------------------------------
\section{تمرین (نارنجی روشن)}
\label{sec:exercise}

\begin{exercise}[تمرین نمونه]
    ثابت کنید که مجموع دو عدد زوج، زوج است.
\end{exercise}

%----------------------------------------
\section{راه حل و اثبات (سبز)}
\label{sec:solution-proof}

\begin{proof}
    چون مربع هر عدد حقیقی نامنفی است، نتیجه حاصل می‌شود.
\end{proof}

\begin{solution}
    فرض کنید $a = 2m$ و $b = 2n$. آنگاه $a + b = 2(m + n)$ که زوج است.
\end{solution}

%----------------------------------------
\section{ارجاعات}
\label{sec:references}

ارجاع به قضیه: \cref{thm:boxes-test}

%----------------------------------------
\section{نتیجه‌گیری}
\label{sec:conclusion}

اگر این سند بدون خطا کامپایل شود:

\begin{itemize}
    \item سه رنگ استفاده می‌شود (آبی، نارنجی، سبز)
    \item \texttt{matincream} به‌عنوان پس‌زمینه استفاده می‌شود
    \item \texttt{fonttitle=\textbackslash normalsize} است
    \item محیط‌های قضیه، لم، تعریف، مثال، نکته، تمرین، راه حل، اثبات کار می‌کنند
    \item \texttt{mb-boxes} با \texttt{mb-theorem} تداخل ندارد
\end{itemize}

\textbf{وضعیت تست:} قبول

\end{document}